\section{Introduction}

Diffusion models bring breakthrough developments in image generation \citep{ldm}, enabling users without any art background to easily create unique and personalized designs, graphics, and illustrations based on specific textual descriptions. Building upon this success, there is a growing interest in extending this concept to video generation \citep{lvdm, ge2023preserve, blattmann2023align, wang2023videocomposer, luo2023videofusion}.
%This application has the potential to revolutionize the entertainment industry by creating concept art for movies or video games based on script descriptions.
%
%However, despite notable advancements in image generation, the field of video generation remains relatively unexplored. This is primarily due to the added complexity introduced by the temporal dimension. 
%
%Generating high-quality videos requires capturing not only the visual content but also the temporal dynamics and smooth transitions between frames.
As targeting for modeling higher dimensional data, video diffusion model demands a notably increased requirement in model capacity and data scale. As a result, current video diffusion models are generally trained on a small number of frames. Consequently, during the inference stage, the quality of the generated video tends to decrease as the length of the video increases due to longer videos are not supervised during the training stage.
%
One straightforward approach is to generate video fragments of the same length as the training videos and then stitch them together, eliminating the training-inference gap. However, this method results in disconnected and incoherent fragments. To address this issue, the fragments can be fused during the denoising process and smoothly connected in the final video \citep{wang2023gen}. 
%
However, the long-distance fragments often have a large content gap to fuse and thus it struggles to maintain the content consistency in the long video. Although some auto-regressive-based methods \citep{villegas2022phenaki} get rid of this problem by progressively generating the next frame, content consistency is still hard to guarantee due to the error accumulation. 

% Diffusion models bring breakthrough developments in image generation  \citep{}. Users without any art background can easily create unique and personalized designs, graphics, and illustrations based on specific textual descriptions. Extending this concept to video generation, it could be used in the entertainment industry to create concept art for movies or video games based on script descriptions.
% %
% To realize this amazing blueprint, recently many diffusion-based video generation models are emerging \citep{}. However, compared to notable advancements in image generation, the realm of video generation remains comparatively unexplored due to the additional complexity introduced by the temporal dimension. 

%
% First of all, those video diffusion models \citep{} are trained on very few frames limited by computing resources. Therefore, during the inference stage, the quality of the video decreases as the length of the generated video increases due to longer videos are not supervised during the training stage. 
% %
% A straightforward idea is to generate fragments with the same length as the training video and then collage them. In this case, there is no connection between the fragments and they are also incoherent. To improve it, those fragments can be fused during the denoising process and be smoothly connected in the final \citep{}. However, the fragments at the beginning and end often have a large content gap to fuse, so it is difficult to maintain the consistency of the content in the synthesized long video. Although some auto-regressive-based methods \citep{} get rid of this problem by progressively generating the next frame, content consistency is still hard to guarantee. 

\if 0
Is there any way to gain some video fragments with similar main subjects and scenarios? To answer this question, we review the inference pipeline of the video diffusion model VideoLDM \citep{blattmann2023align}. 
%
In this model, the denoising process is guided by initializing noise latents and iteratively denoising them using a U-Net\citep{ronneberger2015u} under the given prompt. The final denoised latents are then passed through a pretrained variational autoencoder (VAE) \citep{kingma2013auto} to generate the final video.
%
In VideoLDM\citep{blattmann2023align}, the generated frame depends not only on the initial noise for the current frame but also on the initial noises for all frames. This means that resampling the noise of any frame will significantly influence other frames due to the full interaction facilitated by the temporal attention layers. This makes it challenging to introduce new content while maintaining the main subjects and scenes of the original video.
%
\fi

In VideoLDM\citep{blattmann2023align}, the generated frame depends not only on the initial noise for the current frame but also on the initial noises for all frames. This means that resampling the noise of any frame will significantly influence other frames due to the full interaction facilitated by the temporal attention layers. This makes it challenging to introduce new content while maintaining the main subjects and scenes of the original video.
To address this challenge, we inspect the temporal modeling mechanism of VideoLDM, where the temporal attention module is order-independent, whereas the temporal convolution module is order-dependent. Our experimental observation indicates that the per-frame noises serve as a foundation for determining the overall appearance, while their temporal order influences the content built upon that foundation.
%
Motivated by this, we propose FreeNoise, a tuning-free and time-efficient paradigm to achieve longer video inference. The key idea is to construct a sequence of noise frames with long-range correlation and perform temporal attention over them by the way of window-based fusion. It mainly contains two key designs: \textit{Local Noise Shuffling} and \textit{Window Based Attention Fusion}. By applying the local noise shuffling to a sequence of fixed random noise frames for length extension, we achieve a sequence of noise frames with both internal randomness and long-range correlation. Meanwhile, the window-based attention fusion enables the pre-trained temporal attention modules to process frames of any longer length. Particularly, the overlapped window slicing and merging operation only happens in temporal attention while introducing no computation overhead to other modules of the VideoLDM, which benefits the computational efficiency significantly.

In addition, most video generation models \citep{blattmann2023align, luo2023videofusion, ge2023preserve} only utilize a single-text condition to control the video even when multi-text conditions are given. For instance, the sentence ``A man sleeps on the desk and then reads the book" which contains two stages but only one condition will be reflected in the generated video. This limitation arises from the fact that the training dataset usually contains only a single-text condition. However, in a single-shot scene, the main subject usually involves multiple actions. To address the challenge of generating videos based on multiple prompts without tuning the pretrained models, we propose \textit{Motion Injection}. This approach leverages the characteristics of diffusion models, where different time steps recover varying levels of information (image layout, shapes of the objects, and fine visual details) during the denoising process \citep{patashnik2023localizing,zhang2023prospect}. It gradually injects new motion during the time steps associated with object shapes, following the completion of the previous motion. Importantly, this design does not introduce any additional inference time.

Our contributions are summarized as follows: \textbf{1)} We investigate the temporal modeling mechanism of video diffusion models and identify the influence of initial noises. \textbf{2)} We design a tuning-free paradigm for longer video generation, which outperforms existing state-of-the-art notably in both video quality and computational efficiency. \textbf{3)} We propose an effective motion injection approach that achieves multi-prompt long video generation with decent visual coherence.
